\section{History-complex construction}\label{app:complex}
\subsection*{C.1 The permutohedral route complex}
Let $C=\{N_1,\dots,N_k\}$ be distinct labels and $\mathcal A(m)\subseteq S_k$ the set of admissible arrangements. The boundary complex $\partial P_k$ of the permutohedron~\cite{Ziegler1995} has as faces the ordered set partitions $(B_1|\dots|B_s)$ of $\{1,\dots,k\}$, $s\ge2$; its vertices are the $k!$ arrangements (all blocks singletons) and a face contains an arrangement when the arrangement lists the blocks in order. The \emph{route complex} $K_m$ is the full subcomplex on $\mathcal A(m)$: a face is present iff every arrangement it contains is admissible. The cubical complex $Q_m$ instead has as cells the sets of pairwise non-overlapping legal swaps of one history (\cref{lem:indep}); the two coincide on the $1$-skeleton, differ in higher cells, and the hexagon of three labels is a $2$-cell of $K_m$ but not of $Q_m$.
\subsection*{C.2 Chain complexes and the check $\partial^2=0$}
Both complexes are used with $\mathbb F_2$ coefficients and cellular boundary $\partial\sigma=\sum_{\tau\subset\sigma,\ \dim\tau=\dim\sigma-1}\tau$. For the face poset of $\partial P_k$ (a regular CW sphere) and for the cubical complex every codimension-two incidence is a diamond, so $\partial^2=0$; in addition the program asserts $\partial_{i-1}\partial_i=0$ numerically before each rank computation (\texttt{scripts/topo.py}, function \texttt{betti\_from\_boundaries}), and the Betti numbers are checked against the Euler characteristic of the $f$-vector in every table of this manuscript. Those checks passed in every reported case (for instance $16-20+4=0$ and $1-1+0=0$ for the loop of $\{2,3,8,9\}$; $96-180+98-12=2$ and $1-0+1-0=2$ for the void).
\subsection*{C.3 Implementation}
Cells of $K_m$ are enumerated as ordered set partitions whose linear extensions all lie in $\mathcal A(m)$; ranks over $\mathbb F_2$ are computed by Gaussian elimination on bit-packed rows. For repeated labels no such construction is given (\cref{prob:repeated}). General background on complexes of this type: \cite{Bjorner1995,Kozlov2008}.
